\documentclass[../book.tex]{subfiles}
\begin{document}

% Show the heading here only when compiling this file on its own.
\ifSubfilesClassLoaded{\chapter{EXAMPLE CHAPTER}}{}

This chapter is to provide formatting conventions.

\section*{Preamble that Every Chapter Should Have}

% \noindent\textit{Primary citations/when/who/where course derived}
\noindent\textit{Based on Course ABC from College University (2025-2026) using Book XYZ} \nocite{example2025}
\newline

Discuss motivation, central objects of study, what will be covered, and prerequisites. Prerequisites can vary even within a chapter; e.g. there can be an added real analysis prerequisite beginning in section $n$. Also discuss whether the sections should be read in order. Minor citations can also go here. 

\textbf{Important notes:} 
\begin{enumerate}
    \item Be sure to keep citations up-to-date as you go \cite{example2026}. Citations go to the end of the chapter automatically. Use ``nocite" (like above, in the preamble) to include a citation that isn't in-line.
    \item Keep the index up to date as you go. Here's an example.\index{Example}
\end{enumerate}

\section*{Example Section}

\begin{definition}
    The \textbf{real numbers} are a complete ordered field.
\end{definition}

\begin{lemma}\label{lemma:example}
    This is a lemma.
\end{lemma}

\begin{proposition}
    Propositions are the workhouse result.
    \begin{proof}
        Proofs go inside.
    \end{proof}
\end{proposition}

\begin{theorem}[Named theorems go here, optional]
    Quintics are unsolvable.
\end{theorem}

\begin{theorem}\label{theorem:example-theorem}
    Theorem here that gets quoted later.
\end{theorem}

Here's a standalone proof, apart from its theorem:
\begin{proof}[Proof of Theorem~\ref{theorem:example-theorem}]
    Proofs can reference other results like Lemma~\ref{lemma:example}.
\end{proof}

\begin{remark}
    This is for educational purposes.
\end{remark}

\begin{aside}
    This is for funsies.
\end{aside}

\begin{notation}[Optional text]
    If abbreviations are used, etc.
\end{notation}

\begin{example}\label{example:example-example}
    This is an example. Not all examples have to be in the example environment, but it's useful to reference later.
\end{example}

This is a reference to example~\ref{example:example-example}.

Equations should be like this.
\begin{equation*}
    x+y=z
\end{equation*}
or numbered like
\begin{equation}\label{equation:example}
    x+y=z
\end{equation}
See equation~\ref{equation:example}.
\newline

Some common notation: $\R$, $\Z$, $\Q$, $\C$. Default use $\R^n$. Indicator functions are $\I\{x\in S\}$. Expectation is $\E[X]$; probability is $\Prob(X=x)$.

Arg max is like:
\begin{equation*}
    \hat \beta=\argmax_\beta \Vert y-X\beta\Vert^2
\end{equation*}
Left and right are useful for resizing brackets and parentheses:
\begin{equation*}
    \alpha\left(\frac{x}{y}+\I\left\{X^2\in R\right\}\right)
\end{equation*}
Sets should be written as
\begin{equation*}
    A=\{x\in \Q:x^2\leq 2\}
\end{equation*}
Functions are $f:X\rightarrow Y$, sometimes with the notation $x\mapsto f(x)$.
\newline

\noindent Code can be useful to include with the \texttt{minted} package.
\begin{minted}{python}
    print("Hello world!")
\end{minted}
\noindent Algorithm/pseudocode blocks without specific implementations are nice too.
\begin{algorithm}
\caption{An algorithm with caption}\label{alg:cap}
\begin{algorithmic}
\Require $n \geq 0$
\Ensure $y = x^n$
\State $y \gets 1$
\State $X \gets x$
\State $N \gets n$
\While{$N \neq 0$}
\If{$N$ is even}
    \State $X \gets X \times X$
    \State $N \gets \frac{N}{2}$  \Comment{This is a comment}
\ElsIf{$N$ is odd}
    \State $y \gets y \times X$
    \State $N \gets N - 1$
\EndIf
\EndWhile
\end{algorithmic}
\end{algorithm}

If you want quotation marks, use ``quotes like this". \textit{GRAMMATICAL NOTE:} If you use quotation marks, include the period or comma AFTER the end quotation --- ``hello", ``world" (not ``hello," ``world"). This goes against grammatical conventions, but it aligns better with math and CS writing (think lists of strings). Anyway, we can change later, but at least if we have a set convention from the start it'll be easy to make changes later. 

\printbibliography[
    heading=subbibintoc,
    title={References}
]

\end{document}
